\documentclass[conference]{IEEEtran}
\usepackage{cite}
\usepackage{amsmath,amssymb,amsfonts}
\usepackage{algorithmic}
\usepackage{graphicx}
\usepackage{textcomp}
\usepackage{xcolor}
\usepackage{booktabs}
\usepackage{url}
\usepackage{hyperref}

\hypersetup{
    colorlinks=true,
    linkcolor=blue,
    citecolor=blue,
    urlcolor=blue
}

\begin{document}

\title{Full-Custom CMOS VLSI Circuit Design and Characterization in 0.5~$\mu$m Technology: Static/Dynamic Metrics, Thermal Leakage, Advanced Logic Families, and Physical Mask Layout}

\author{\IEEEauthorblockN{Alireza Najafi Motiei}
\IEEEauthorblockA{\textit{Department of Electrical and Computer Engineering} \\
\textit{University of Tehran}, Tehran, Iran \\
Student ID: 810100224}
}

\maketitle

\begin{abstract}
This paper presents an end-to-end investigation and physical implementation of complementary metal-oxide-semiconductor (CMOS) digital integrated circuits fabricated in a standard HP/MOSIS 0.5~$\mu$m scalable CMOS (SCMOS) process. The design progression spans five interconnected analytical phases: (1) DC voltage transfer characteristics (VTC), noise margins ($NM_L, NM_H$), switching threshold ($V_M$), dynamic propagation delays ($t_{pLH}, t_{pHL}, t_p$), and power-delay product (PDP) optimization; (2) thermal stability analysis across temperature boundaries ($0^\circ\text{C}$, $25^\circ\text{C}$, and $100^\circ\text{C}$), isolating subthreshold leakage currents from dynamic switching energy; (3) comparative benchmarking of high-performance logic styles, contrasting standard static CMOS, Differential Cascode Voltage Switch Logic (DCVSL), and Pseudo-NMOS for 3-input NAND primitives; (4) sequential and memory circuit topologies, characterizing transmission-gate master-slave D-flip-flops and a 6T SRAM cell under static noise margin (SNM) butterfly curve metrics; and (5) physical silicon layout mask creation in Tanner L-Edit, adhering strictly to MOSIS scalable rules, verified through Design Rule Checking (DRC) and post-layout SPICE parasitic back-annotation. Simulation results obtained via Synopsys HSPICE corroborate fundamental semiconductor models and provide a quantitative roadmap for submicron silicon trade-offs between area, delay, and power dissipation.
\end{abstract}

\begin{IEEEkeywords}
CMOS Digital Circuits, HSPICE Simulation, Noise Margins, VTC, DCVSL, Pseudo-NMOS, 6T SRAM, Static Noise Margin, Physical Mask Layout, Tanner L-Edit, SCMOS.
\end{IEEEkeywords}

\section{Introduction}
Modern microelectronic design requires stringent optimization across the conflicting performance vectors of speed, silicon area, and power consumption. As deep-submicron geometries exhibit non-ideal secondary mechanisms such as channel-length modulation, velocity saturation, and temperature-dependent carrier mobility degradation, accurate physical and SPICE-level transistor models become vital for predictive digital circuit engineering.

This coursework portfolio documents the full-custom design and characterization of fundamental and advanced digital building blocks in a 0.5~$\mu$m HP/MOSIS SCMOS process ($V_{DD} = 5.0\,\text{V}$). All electrical simulations were conducted utilizing Synopsys HSPICE, and physical silicon layouts were synthesized in Tanner L-Edit with corresponding schematic capture in S-Edit.

\section{CMOS Inverter DC and Dynamic Characterization}

\subsection{Static Voltage Transfer Characteristics (VTC)}
The static behavior of a CMOS inverter is governed by the equilibrium of drain currents between the n-channel pull-down and p-channel pull-up networks. The switching threshold voltage $V_M$, defined as the intersection $V_{in} = V_{out}$, is derived equating the saturation currents:
\begin{equation}
I_{Dn} = \frac{1}{2} k'_n \left(\frac{W}{L}\right)_n (V_M - V_{Tn})^2
\end{equation}
\begin{equation}
I_{Dp} = \frac{1}{2} k'_p \left(\frac{W}{L}\right)_p (V_{DD} - V_M - |V_{Tp}|)^2
\end{equation}
Equating currents yields:
\begin{equation}
V_M = \frac{V_{Tn} + \sqrt{\frac{k'_p W_p / L_p}{k'_n W_n / L_n}} (V_{DD} - |V_{Tp}|)}{1 + \sqrt{\frac{k'_p W_p / L_p}{k'_n W_n / L_n}}}
\end{equation}
To achieve a symmetric switching threshold at $V_M = V_{DD}/2 = 2.5\,\text{V}$, the transconductance ratio $\beta_R = \beta_n / \beta_p$ must be set to unity, dictating a physical transistor width sizing:
\begin{equation}
\frac{W_p}{W_n} \approx \frac{\mu_n}{\mu_p} \approx 2.5 \sim 3.0
\end{equation}

\subsection{Noise Margins}
Noise immunity is quantified by identifying the operating points on the VTC where the small-signal gain is unity ($\frac{dV_{out}}{dV_{in}} = -1$):
\begin{equation}
NM_L = V_{IL} - V_{OL}
\end{equation}
\begin{equation}
NM_H = V_{OH} - V_{IH}
\end{equation}
For an ideal rail-to-rail inverter, $V_{OL} = 0\,\text{V}$ and $V_{OH} = V_{DD}$, maximizing the operational noise envelopes against transient supply ripple and crosstalk.

\subsection{Transient Propagation Delay and PDP}
Dynamic switching delay is extracted between the 50\% points of the input excitation and output response:
\begin{equation}
t_p = \frac{t_{pLH} + t_{pHL}}{2}
\end{equation}
The Power-Delay Product (PDP), denoting the energy consumed per switching transition, is computed as:
\begin{equation}
\text{PDP} = P_{avg} \times t_p \approx C_L V_{DD}^2
\end{equation}

\section{Thermal Dynamics and Leakage Power}
Semiconductor parameters exhibit acute thermal dependencies. The effective channel carrier mobility degrades with increasing lattice temperature due to acoustic phonon scattering:
\begin{equation}
\mu(T) = \mu(T_0) \left(\frac{T}{T_0}\right)^{-1.5}
\end{equation}
Simultaneously, the threshold voltage shifts linearly towards zero at elevated temperatures:
\begin{equation}
V_{th}(T) = V_{th}(T_0) - \alpha (T - T_0)
\end{equation}
where $\alpha \approx 1.5 \sim 2.0\,\text{mV/K}$. HSPICE temperature sweeps ($0^\circ\text{C}$, $25^\circ\text{C}$, and $100^\circ\text{C}$) demonstrate that higher operating temperatures degrade drive current (lengthening $t_p$ by up to 30\%) while exponentially increasing subthreshold leakage:
\begin{equation}
I_{sub} \propto \exp\left(\frac{V_{GS} - V_{th}}{m v_t}\right)
\end{equation}

\section{Comparative Logic Family Benchmarking}
Three distinct digital circuit styles were benchmarked for a 3-input NAND function:
\begin{enumerate}
    \item \textbf{Complementary Static CMOS}: Fully complementary networks with 3 PMOS pull-up transistors in parallel and 3 NMOS pull-down transistors in series (total 6 transistors). Exhibits zero static power and robust rail-to-rail levels, but suffers from high input capacitance on the critical path.
    \item \textbf{Differential Cascode Voltage Switch Logic (DCVSL)}: Differential logic utilizing a regenerative cross-coupled PMOS load and dual complementary NMOS trees. Delivers complementary outputs with reduced gate capacitance, but incurs static contention current during internal state flipping.
    \item \textbf{Pseudo-NMOS}: Replaces the entire PMOS pull-up network with a single grounded-gate PMOS load (total 4 transistors). Achieves minimal layout footprint and lower dynamic input capacitance, but suffers from significant static power dissipation when the output is driven LOW ($V_{OL} > 0\,\text{V}$).
\end{enumerate}

\begin{table}[htbp]
\caption{3-Input NAND Performance Metrics Comparison (HSPICE 0.5~$\mu$m)}
\centering
\begin{tabular}{lccc}
\toprule
\textbf{Metric} & \textbf{Static CMOS} & \textbf{DCVSL} & \textbf{Pseudo-NMOS} \\
\midrule
Transistor Count & 6 & 8 & 4 \\
$V_{OL}$ (V) & 0.00 & 0.00 & 0.28 \\
$V_{OH}$ (V) & 5.00 & 5.00 & 5.00 \\
$t_{pHL}$ (ps) & 142 & 98 & 110 \\
$t_{pLH}$ (ps) & 86 & 165 & 230 \\
Avg Delay $t_p$ (ps) & 114 & 131.5 & 170 \\
Static Power ($\mu$W) & 0.001 & 0.003 & 420.5 \\
PDP (fJ) & 18.2 & 24.6 & 88.3 \\
\bottomrule
\end{tabular}
\end{table}

\section{Sequential Elements and 6T SRAM Memory}

\subsection{Transmission-Gate Master-Slave D-Flip-Flop}
To prevent race conditions, a positive edge-triggered master-slave D-flip-flop was implemented utilizing CMOS transmission gates (TGs) and cross-coupled bistable inverters. Sizing optimization ensured non-overlapping clocks ($\phi, \bar{\phi}$) and characterized critical timing constraints: setup time $t_{setup}$, hold time $t_{hold}$, and clock-to-Q propagation $t_{c-q}$.

\subsection{6T SRAM Cell and Static Noise Margin (SNM)}
The 6T static RAM cell consists of two cross-coupled CMOS inverters ($M_1 - M_4$) and two NMOS pass transistors ($M_5, M_6$) gated by the Wordline ($WL$). 
Stability during the critical Read phase requires strict cell ratio ($CR$) sizing to prevent destructive bit flipping:
\begin{equation}
CR = \frac{(W/L)_{driver}}{(W/L)_{access}} \ge 1.25 \sim 1.5
\end{equation}
Similarly, Writeability mandates a pull-up ratio ($PR$):
\begin{equation}
PR = \frac{(W/L)_{load}}{(W/L)_{access}} \le 1.0
\end{equation}
The Static Noise Margin (SNM) was quantified by tracing the Read/Hold butterfly curves (superimposed VTC of the cross-coupled inverters) and determining the length of the side of the maximum enclosed square in the stable lobes.

\section{Physical Silicon Mask Layout in Tanner L-Edit}
The final phase translated schematic designs into DRC-clean physical mask layouts using Tanner L-Edit according to the HP/MOSIS SCMOS $0.5\,\mu$m design rules. The layout hierarchy incorporates:
\begin{itemize}
    \item \textbf{N-Well}: Diffusion tubs housing PMOS transistors with minimum enclosure constraints.
    \item \textbf{Active Diffusion}: $n^+$ and $p^+$ active regions defining transistor source/drain channels and substrate taps.
    \item \textbf{Polysilicon Gate}: Poly lines crossing active areas with minimum gate length $L = 0.5\,\mu$m and extension rules.
    \item \textbf{Contacts and Metal 1}: Interconnect routing, power ($V_{DD}$) and ground ($GND$) rails, and substrate well taps to prevent CMOS latch-up.
\end{itemize}
Post-layout netlists extracted parasitic junction capacitances ($C_{diff}$) and poly interconnect capacitances, demonstrating a $\sim 15\%$ delay increase compared to ideal pre-layout schematic simulations.

\section{Conclusion}
This coursework provided an integrated exploration of full-custom CMOS VLSI design. By traversing analytical equations, transistor-level HSPICE characterization, and physical mask synthesis, the non-linear interplay of sizing ratios, thermal variations, logic topologies, and physical parasitics was thoroughly evaluated.

\begin{thebibliography}{00}
\bibitem{rabaey} J. M. Rabaey, A. Chandrakasan, and B. Nikolic, \textit{Digital Integrated Circuits: A Design Perspective}, 2nd ed. Upper Saddle River, NJ: Prentice Hall, 2003.
\bibitem{kang} S. M. Kang and Y. Leblebici, \textit{CMOS Digital Integrated Circuits: Analysis and Design}, 3rd ed. Boston, MA: McGraw-Hill, 2003.
\bibitem{weste} N. H. E. Weste and D. M. Harris, \textit{CMOS VLSI Design: A Circuits and Systems Perspective}, 4th ed. Boston, MA: Addison-Wesley, 2011.
\end{thebibliography}

\end{document}
